\documentclass[11pt]{article}

\usepackage{amsmath,amssymb,graphicx,booktabs,geometry,hyperref,natbib}

\geometry{margin=1in}
\hypersetup{
  colorlinks=true,
  linkcolor=blue,
  urlcolor=blue,
  citecolor=blue
}

\title{A Homogeneous Compensation Mechanism for Indirect Evolutionary Rescue in a Roy-Style Predator--Prey Model}
\author{Mathbio research notes}
\date{\today}

\begin{document}

\maketitle

\begin{abstract}
Indirect evolutionary rescue can occur when evolutionary change in one species prevents extinction of another. We analyze this mechanism in a Roy-style predator--prey model with evolving prey defense, comparing a well-mixed eco-evolutionary ODE with a spatial PDE using the same local reactions. The central mechanism is a homogeneous compensation branch: as predator mortality stress increases, the equilibrium prey defense frequency decreases, increasing predator conversion opportunity and preserving a positive predator equilibrium over an interior stress interval. For the linear defense trade-off ODE, this branch follows from explicit feasibility conditions and is locally stable at the tested stresses by the Routh--Hurwitz inequalities for the cubic characteristic polynomial. In the spatial PDE, the same homogeneous branch is linearly stable to tested Neumann modes. Targeted finite-amplitude heterogeneous perturbations produce no persistent spatial-pattern-mediated rescue in the tested setup; finite-horizon basin-boundary excursions resolve to matched homogeneous controls. A controlled nonlinear trade-off extension with concave, convex, and mixed endpoint-preserving shapes supports compensation beyond the linear case, while remaining local and parameter-sensitive. The evidence supports homogeneous reaction-level compensation, not spatial-pattern-mediated rescue, in the tested model.
\end{abstract}

\section{Introduction}

Evolutionary rescue describes persistence after environmental deterioration through adaptive change \citep{GomulkiewiczHolt1995,BellGonzalez2009,Bell2017}. In predator--prey systems, rescue can also be indirect: evolution in prey may alter predator persistence. In the setting studied here, prey defense reduces predation but carries ecological costs. When predator mortality stress increases, predator pressure declines, selection can shift prey toward lower defense, and the lower-defense prey state can increase predator conversion opportunity. This feedback is the qualitative mechanism of indirect evolutionary rescue described by \citet{YamamichiMiner2015}.

Prey defense evolution is a natural target for this question because the defense trait changes both prey ecological performance and predator energetic gain. Rapid eco-evolutionary feedbacks in predator--prey systems can alter ecological dynamics on short timescales \citep{YoshidaEtAl2003}. The exact form of the defense trade-off is also important: nonlinear or asymmetric trade-offs can change predator--prey eco-evolutionary dynamics \citep{KasadaYamamichiYoshida2014}.

Spatial structure is relevant because local dispersal and spatial heterogeneity can change predator--prey overlap, delay extinction, and alter basin entry. A scalar persistence threshold is an attractive first diagnostic: one compares the mortality stress at which predators persist or fail in the ODE and PDE. That threshold framing can fail, however, when long transients, continuation dependence, or multiple reachable outcomes occur at the same stress. In that case, the correct object is not a single threshold but the mechanism organizing the reachable basins.

The mechanism supported by the present analysis is not spatial-pattern-mediated rescue. Fixed-defense spatial patterning did not produce robust predator rescue. The well-mixed eco-evolutionary ODE did support indirect evolutionary rescue. Subsequent ODE--PDE comparison showed that the basin structure observed in the PDE is primarily inherited from homogeneous reaction dynamics rather than generated by persistent spatial patterning. The important mechanism is therefore a homogeneous compensation branch, with spatial analyses used to test whether the PDE destabilizes or qualitatively alters that branch.

We ask whether the rescue mechanism is controlled by spatial patterning or by homogeneous eco-evolutionary reaction dynamics. We then derive the compensation branch responsible for rescue, test its local stability, test whether the spatial PDE destabilizes it, and examine whether the mechanism persists under controlled nonlinear trade-off shapes.

The claim hierarchy is intentionally conservative. Strong claims are restricted to the tested ODE branch, its feasibility and local stability, the tested PDE spatial-mode stability calculation, targeted non-homogeneous perturbation outcomes, and the controlled nonlinear shape grid. Broader claims about global stability, all trade-off forms, all diffusion coefficients, or all heterogeneous initial conditions are not made.

\section{Model}

Let \(n\) be total prey density, \(w\) predator density, \(q\) prey defense frequency, and
\[
z=\kappa^{-1}-n-w
\]
be free space. For the linear trade-off model,
\[
r(q)=r_u(1-q)+r_vq,
\quad
a(q)=a_u(1-q)+a_vq,
\quad
b(q)=b_u(1-q)+b_vq.
\]
The well-mixed ODE is
\begin{align}
\frac{dn}{dt} &= n[r(q)z-\xi-a(q)w],\\
\frac{dw}{dt} &= w[b(q)nz-(m+s)-\mu w],\\
\frac{dq}{dt} &= \nu q(1-q)[(r_v-r_u)z-(a_v-a_u)w].
\end{align}
Here \(s\) is additional predator mortality stress. Table~\ref{tab:parameters} lists the values used in the reported calculations.

The spatial PDE uses the same local reaction terms:
\begin{align}
\partial_t n &= D_n\nabla^2n+n[r(q)z-\xi-a(q)w],\\
\partial_t w &= D_w\nabla^2w+w[b(q)nz-(m+s)-\mu w],\\
\partial_t q &= D_q\nabla^2q+\nu q(1-q)[(r_v-r_u)z-(a_v-a_u)w].
\end{align}
The PDE tests use zero-flux Neumann boundaries.

\begin{table}[htbp]
  \centering
  \caption{Parameter values used in the compensation-branch and spatial-stability analyses. These values define the tested mechanistic case study; they are not empirical estimates.}
  \label{tab:parameters}
  \begin{tabular}{llll}
    \toprule
    Symbol & Value & Symbol & Value \\
    \midrule
    \(r_u\) & 1 & \(r_v\) & 0.65 \\
    \(a_u\) & 1 & \(a_v\) & 0.35 \\
    \(b_u\) & 0.08 & \(b_v\) & 0.02 \\
    \(\kappa\) & 0.15 & \(\xi\) & 0.55 \\
    \(m\) & 0.1 & \(\mu\) & 0.2 \\
    \(\nu\) & 0.05 & \(D_n\) & 0.01 \\
    \(D_w\) & 0.01 & \(D_q\) & 0.005 \\
    \(L_x\) & 20 & \(L_y\) & 20 \\
    \bottomrule
  \end{tabular}
\end{table}

For the controlled nonlinear trade-off extension, endpoint values are preserved but linear interpolation is replaced by shape functions:
\[
r(q)=r_u+(r_v-r_u)f_r(q),
\]
\[
a(q)=a_u+(a_v-a_u)f_a(q),
\]
\[
b(q)=b_u+(b_v-b_u)f_b(q),
\]
with \(f_\gamma(q)=q^\gamma\). This nonlinear form is used only for controlled extension tests, not to replace the baseline linear model.

\section{Methods and numerical protocols}

The analysis combines targeted numerical integration with analytic or semi-analytic calculations. ODE trajectories were integrated from specified initial means and classified from late-time behavior of \(w(t)\), \(q(t)\), relative change between terminal windows, and residual-style steady-state diagnostics. Predator-persistent outcomes have positive predator density with sufficiently small terminal change; extinct outcomes have predator density near zero; transient outcomes do not settle clearly into either class within the tested horizon. The same classification vocabulary is used for matched PDE spatial means, with additional spatial coefficient-of-variation diagnostics for pattern strength.

The homogeneous compensation branch was computed from the interior equilibrium equations. For the linear trade-off model this gives the closed-form branch \(z^*,w^*,n^*,q^*(s)\). Numerical equilibria from multiple initial guesses were used only as a consistency check against the analytic branch. Local stability was evaluated with the analytic reaction Jacobian. For the cubic characteristic polynomial \(\lambda^3+A_1\lambda^2+A_2\lambda+A_3=0\), the Routh--Hurwitz inequalities \(A_1>0\), \(A_2>0\), \(A_3>0\), and \(A_1A_2>A_3\) determined local stability.

PDE spatial-mode stability was evaluated by treating each homogeneous branch state \(U^*\) as a spatially constant PDE steady state and computing eigenvalues of \(J_F(U^*)-\lambda_{ij}D\) for Neumann modes. The zero mode recovers the ODE stability calculation; nonzero modes test infinitesimal spatial perturbations. Finite-amplitude PDE tests used localized predator patches, localized defense patches, sinusoidal modes, random heterogeneity, and basin-boundary heterogeneity, each compared with a matched homogeneous control. Basin-changing heterogeneous cases were followed to longer horizons to distinguish persistent basin changes from finite-horizon transients.

The nonlinear trade-off extension used the endpoint-preserving shape family \(f_\gamma(q)=q^\gamma\) with \(\gamma\in\{0.5,1,2\}\) for the growth, attack, and conversion trade-offs. This is a controlled local shape grid, not an exhaustive parameter scan. Across all numerical protocols, tail-window and residual checks were used to avoid interpreting unresolved long transients as steady outcomes.

The main text uses a reduced figure set focused on mechanism, stability, spatial perturbation response, and nonlinear extension. Additional diagnostic panels from the underlying analyses are treated as supporting evidence rather than main-text figures.

\section{Homogeneous compensation branch}

Define
\[
\Delta r=r_v-r_u,\quad \Delta a=a_v-a_u,\quad \Delta b=b_v-b_u.
\]
For an interior equilibrium with \(n>0\), \(w>0\), and \(0<q<1\), the nontrivial equilibrium conditions are
\[
r(q)z-\xi-a(q)w=0,
\]
\[
b(q)nz-(m+s)-\mu w=0,
\]
\[
\Delta r z-\Delta a w=0.
\]
If
\[
c=\frac{\Delta r}{\Delta a},
\]
then the selection-gradient condition gives
\[
w^*=cz^*.
\]
Because the trade-offs are linear, \(r(q)-ca(q)=r_u-ca_u\), so the prey-balance equation gives
\[
z^*=\frac{\xi}{r_u-ca_u},
\]
\[
w^*=cz^*,
\]
\[
n^*=\kappa^{-1}-z^*-w^*.
\]
The predator-balance equation then determines the defense frequency required at stress \(s\):
\[
q^*(s)=\frac{(m+s+\mu w^*)/(n^*z^*)-b_u}{b_v-b_u}.
\]

This branch explains rescue as compensation: increasing predator mortality stress is offset by decreasing equilibrium prey defense frequency, which increases predator conversion opportunity and keeps a positive predator equilibrium feasible. In the current parameterization, the branch has approximately \(n^*\approx 4.8333\) and \(w^*\approx 0.6417\), while \(q^*(s)\) decreases with stress.

\begin{figure}[htbp]
  \centering
  \begin{minipage}{0.49\linewidth}
    \centering
    \includegraphics[width=\linewidth]{figures/roy_evo_spatial/report/fig33_compensation_branch_current.png}
  \end{minipage}
  \hfill
  \begin{minipage}{0.49\linewidth}
    \centering
    \includegraphics[width=\linewidth]{figures/roy_evo_spatial/report/fig40_compensation_stress_interval.png}
  \end{minipage}
  \caption{Compensation branch and stress interval. Left: analytic branch for the linear trade-off parameterization, compared with numerical stable persistent equilibria. Right: stress interval where \(0<q^*(s)<1\). Together these panels show that increasing mortality stress is balanced by lower prey defense over the tested branch interval. The allowed conclusion is branch feasibility and analytic/numerical agreement in the tested parameterization, not global stability.}
  \label{fig:branch-interval}
\end{figure}

\section{Existence and stress interval}

The interior branch exists only under explicit feasibility conditions:
\[
c>0,
\]
\[
r_u-ca_u>0,
\]
\[
z^*>0,\quad w^*>0,\quad n^*>0,
\]
\[
0<q^*(s)<1.
\]
Solving the branch equation for stress gives
\[
s(q)=n^*z^*[b_u+\Delta bq]-m-\mu w^*.
\]
Therefore the open stress interval for an interior branch is the interval between the values at \(q=0\) and \(q=1\). In the current parameterization this is approximately
\[
s\in(-0.1131,0.2324).
\]
The target stresses used in the mechanism tests lie inside this interval. These conditions make the branch a conditional mathematical statement for the linear trade-off ODE rather than a purely numerical observation.

\section{Local stability of the compensation branch}

Local stability of the homogeneous branch was evaluated with the analytic Jacobian of the ODE. For a three-dimensional system, the characteristic polynomial can be written
\[
\lambda^3+A_1\lambda^2+A_2\lambda+A_3=0.
\]
The cubic Routh--Hurwitz conditions for local stability are
\[
A_1>0,\quad A_2>0,\quad A_3>0,\quad A_1A_2>A_3.
\]
For the current parameterization, these inequalities hold at all target stresses tested along the branch and agree with direct eigenvalue stability. The smallest recorded Routh--Hurwitz margin \(A_1A_2-A_3\) across the target stresses was positive. This supports local stability of the homogeneous compensation branch in the tested stress interval.

This is a local stability result. It does not prove global attraction or determine all basin boundaries.

\begin{figure}[htbp]
  \centering
  \includegraphics[width=0.90\linewidth]{figures/roy_evo_spatial/report/fig42_routh_hurwitz_terms.png}
  \caption{Routh--Hurwitz coefficient terms along the compensation branch. The plotted quantities are \(A_1\), \(A_2\), \(A_3\), and \(A_1A_2-A_3\). Their positivity supports local stability at the tested stresses, but does not establish global basin attraction.}
  \label{fig:rh-terms}
\end{figure}

\section{Spatial PDE stability and non-homogeneous perturbations}

A homogeneous ODE equilibrium is also a homogeneous PDE steady state because diffusion terms vanish for spatially constant fields. For Neumann spatial modes on a rectangular domain, perturbations evolve according to
\[
J_F(U^*)-\lambda_{ij}D,
\]
where \(J_F(U^*)\) is the reaction Jacobian at the homogeneous branch state, \(D=\mathrm{diag}(D_n,D_w,D_q)\), and
\[
\lambda_{ij}=\left(\frac{i\pi}{L_x}\right)^2+\left(\frac{j\pi}{L_y}\right)^2.
\]
The zero mode is the ODE stability problem. Nonzero modes test whether spatial perturbations grow.

For the current diffusion coefficients, all tested nonzero Neumann modes were stable at the target branch stresses. A controlled diffusion-ratio linear grid also found no spatial instability in the tested ratios. Thus the PDE does not destabilize the homogeneous branch through the tested infinitesimal spatial modes.

\begin{figure}[htbp]
  \centering
  \begin{minipage}{0.49\linewidth}
    \centering
    \includegraphics[width=\linewidth]{figures/roy_evo_spatial/report/fig45_pde_mode_growth_rates.png}
  \end{minipage}
  \hfill
  \begin{minipage}{0.49\linewidth}
    \centering
    \includegraphics[width=\linewidth]{figures/roy_evo_spatial/report/fig57_long_horizon_decision.png}
  \end{minipage}
  \caption{Spatial-mode and finite-amplitude PDE tests. Left: maximum real growth rates for tested Neumann modes, computed from \(J_F(U^*)-\lambda_{ij}D\). Right: long-horizon follow-up of finite-amplitude heterogeneous basin-changing cases. These diagnostics support linear spatial stability and show that followed finite-horizon basin changes resolved without persistent spatial patterning.}
  \label{fig:pde-spatial-nonlinear}
\end{figure}

Linear spatial stability is not enough to rule out finite-amplitude heterogeneous effects. Targeted finite-amplitude PDE tests therefore used localized predator patches, localized defense patches, sinusoidal perturbations, smooth random heterogeneity, and basin-boundary heterogeneity. These tests compared heterogeneous fields with matched homogeneous controls. Initial finite-horizon basin-label changes occurred for localized defense perturbations near basin-boundary mean states, but these changes were transient-label effects without persistent spatial patterning. Longer-horizon follow-up showed that the basin-changing cases resolved to the matched homogeneous-control basin, with final spatial coefficients of variation below the persistence threshold.

The spatial interpretation is therefore specific: in the tested PDE, spatial heterogeneity can create transient excursions near basin boundaries, but it did not generate a persistent spatial-pattern-mediated rescue mechanism.

\section{Extension to nonlinear trade-off shapes}

For nonlinear trade-offs, the branch is most naturally parameterized by \(q\). When derivatives are defined and \(a'(q)\ne0\), the selection-gradient condition gives
\[
c(q)=\frac{r'(q)}{a'(q)}.
\]
The corresponding branch geometry is
\[
z(q)=\frac{\xi}{r(q)-a(q)c(q)},
\]
\[
w(q)=c(q)z(q),
\]
\[
n(q)=\kappa^{-1}-z(q)-w(q),
\]
and predator balance maps the branch to stress:
\[
s(q)=b(q)n(q)z(q)-m-\mu w(q).
\]
The linear branch is recovered when all shape exponents are one.

The controlled nonlinear extension used \(\gamma_r,\gamma_a,\gamma_b\in\{0.5,1,2\}\), representing concave, linear, and convex endpoint-preserving shapes. This grid is small and structured; it is not a global trade-off scan. Within that grid, the generalized formulation recovered the linear branch, and selected concave, convex, and mixed nonlinear shapes supported feasible locally stable compensation branches. The shape-grid summary recorded 11 robust compensation shapes, 11 partial compensation shapes, 4 no-compensation shapes, and 1 unresolved shape. Targeted PDE spatial-mode and non-homogeneous tests for selected stable nonlinear branches did not support persistent spatial-pattern-mediated rescue.

The nonlinear result suggests that compensation is not only a linear-trade-off artifact within this controlled shape grid. It also emphasizes parameter sensitivity: nonlinear shape can change branch feasibility, stability, and stress interval length.

\begin{figure}[htbp]
  \centering
  \begin{minipage}{0.53\linewidth}
    \centering
    \includegraphics[width=\linewidth]{figures/roy_evo_spatial/report/fig59_nonlinear_branch_curves.png}
  \end{minipage}
  \hfill
  \begin{minipage}{0.43\linewidth}
    \centering
    \includegraphics[width=\linewidth]{figures/roy_evo_spatial/report/fig64_nonlinear_tradeoff_final_decision.png}
  \end{minipage}
  \caption{Controlled nonlinear trade-off extension. Left: selected nonlinear branch curves plotting stress as a function of branch defense frequency. Right: summary of branch recovery, shape-grid outcomes, and targeted PDE checks. These panels support compensation in selected nonlinear shapes while showing that the result remains local and parameter-sensitive.}
  \label{fig:nonlinear-extension}
\end{figure}

\section{Biological interpretation}

Higher predator mortality directly reduces predator persistence. In this model, prey defense evolution can compensate because mortality stress reduces predator pressure, and lower predator pressure favors lower prey defense. Lower defense increases predator conversion opportunity through \(b(q)\), which can restore predator growth balance at a positive predator equilibrium. This is indirect evolutionary rescue: the predator is rescued not by predator adaptation, but by prey evolution.

The biological interpretation is therefore a compensation story rather than a spatial patterning story. Defended prey are less favorable to predators in the conversion term, so a shift toward lower defense makes prey more exploitable by the predator. Under the tested trade-off values, that evolutionary shift is sufficient to offset part of the imposed predator mortality stress. The predator persists indirectly because the prey population evolves toward a state that improves predator growth balance.

The spatial PDE does not create the rescue mechanism in the tested setup. Instead, the PDE preserves a homogeneous reaction-level branch that is locally stable to tested spatial modes. Spatial heterogeneity can matter transiently near basin boundaries, but the targeted tests did not produce persistent spatial patterning or spatial-pattern-mediated rescue.

Several assumptions are stylized. The defense trait is represented by a single frequency \(q\), trade-offs are low-dimensional functions, diffusion is isotropic and constant, and the environmental stress enters as an added predator mortality term. These choices make the mechanism interpretable but not biologically calibrated. Empirical prediction would require species-specific estimates of defense costs, predator conversion rates, movement scales, mortality stress, and the timescale of prey evolutionary response.

\begin{figure}[htbp]
  \centering
  \includegraphics[width=0.88\linewidth]{figures/roy_evo_spatial/report/fig58_final_mechanism_synthesis.png}
  \caption{Final mechanism synthesis. The diagram summarizes the permitted mechanism claim: predator mortality stress shifts homogeneous eco-evolutionary balance, lower prey defense increases predator conversion opportunity, and the spatial PDE preserves this reaction-level mechanism in the tested setup. It does not imply spatial-pattern-mediated rescue.}
  \label{fig:final-synthesis}
\end{figure}

\section{Limitations}

The linear compensation branch result is exact for the linear trade-off ODE under the stated feasibility conditions, but it is not automatically a theorem for all trade-off forms. The nonlinear trade-off extension is controlled and local, using only a structured exponent grid. The PDE tests are targeted rather than exhaustive. There is no global basin proof, no proof of global stability, no theorem for all heterogeneous initial conditions, and no full bifurcation analysis.

The model is not biologically calibrated, and the parameters are not fitted to empirical data. Parameters are used to test mechanism, not to estimate a particular predator--prey system. The results should therefore be interpreted as a mechanistic case study: they support homogeneous reaction-level compensation in the tested Roy-style model, not a universal claim about spatial evolutionary rescue.

\section{Claims and scope}

The established result in this parameterization is that the well-mixed eco-evolutionary ODE has an interior compensation branch, that the branch is feasible over an explicit stress interval, and that the branch is locally stable at the tested stresses by Routh--Hurwitz criteria. For the spatial PDE, the supported result is that the same homogeneous branch is linearly stable to tested Neumann modes and that targeted finite-amplitude heterogeneous perturbations do not produce persistent spatial-pattern-mediated rescue.

The controlled extension supports the presence of compensation branches beyond the linear trade-off form for selected concave, convex, and mixed endpoint-preserving shapes. This does not establish global robustness over all nonlinear trade-offs. The manuscript does not claim global stability, a complete basin theorem, empirical calibration, spatial-pattern-mediated rescue, or exhaustive coverage of heterogeneous initial conditions.

\section{Conclusion}

\begin{enumerate}
  \item The tested ODE supports indirect evolutionary rescue through prey defense evolution.
  \item The mechanism is a homogeneous compensation branch.
  \item Branch existence follows from explicit feasibility conditions.
  \item The current branch is locally stable by Routh--Hurwitz criteria.
  \item The spatial PDE branch is linearly stable to tested Neumann spatial modes.
  \item Targeted non-homogeneous perturbations do not produce persistent spatial-pattern-mediated rescue.
  \item Controlled nonlinear trade-off tests show that the mechanism extends beyond the linear case but remains parameter-sensitive.
\end{enumerate}

The current evidence supports homogeneous reaction-level compensation, not spatial-pattern-mediated rescue, in the tested model.

\appendix

\section{Derivation of the compensation branch}

At an interior equilibrium, the prey equation, predator equation, and selection-gradient equation are all zero after division by the positive factors \(n\), \(w\), and \(\nu q(1-q)\). The selection-gradient equation gives
\[
\Delta r z-\Delta a w=0.
\]
If \(c=\Delta r/\Delta a>0\), this yields \(w=cz\). Substituting into prey balance,
\[
r(q)z-\xi-a(q)cz=0.
\]
For linear trade-offs,
\[
r(q)-ca(q)=r_u-ca_u,
\]
so
\[
z^*=\frac{\xi}{r_u-ca_u},\quad w^*=cz^*,\quad n^*=\kappa^{-1}-z^*-w^*.
\]
Finally, predator balance gives
\[
b(q^*)=\frac{m+s+\mu w^*}{n^*z^*},
\]
and because \(b(q)=b_u+\Delta bq\),
\[
q^*(s)=\frac{(m+s+\mu w^*)/(n^*z^*)-b_u}{\Delta b}.
\]

\section{Jacobian and Routh--Hurwitz coefficients}

Write
\[
F_n=nR,\quad R=r(q)z-\xi-a(q)w,
\]
\[
F_w=wP,\quad P=b(q)nz-(m+s)-\mu w,
\]
\[
F_q=\nu q(1-q)G,\quad G=\Delta r z-\Delta a w.
\]
At the branch, \(R=P=G=0\). The analytic Jacobian was evaluated at the branch and summarized by the characteristic polynomial
\[
\det(\lambda I-J)=\lambda^3+A_1\lambda^2+A_2\lambda+A_3.
\]
The coefficients were computed as
\[
A_1=-\mathrm{tr}(J),
\]
\[
A_2=\frac{1}{2}\left[(\mathrm{tr}J)^2-\mathrm{tr}(J^2)\right],
\]
\[
A_3=-\det(J).
\]
The Routh--Hurwitz stability inequalities are \(A_1>0\), \(A_2>0\), \(A_3>0\), and \(A_1A_2>A_3\).

\section{Numerical protocols and classification rules}

The analyses reported here consolidate existing numerical outputs. They do not introduce new simulations or parameter scans. ODE/PDE comparisons used matched initial means. Spatial-mode tests used Neumann eigenvalues and the modal matrix \(J_F(U^*)-\lambda_{ij}D\). Non-homogeneous PDE tests compared each heterogeneous run with a matched homogeneous control at the same stress and baseline state. Basin-label changes were interpreted only after long-horizon follow-up when available. Persistent spatial patterning was assessed with final spatial coefficient-of-variation thresholds from the corresponding research notes.

\bibliographystyle{plainnat}
\bibliography{manuscript/references}

\end{document}
